\chapter{Dasar-dasar Prompt Engineering}
\label{chap:prompting}

Prompt Engineering sering disebut sebagai "bahasa pemrograman baru" (The New Coding). Jika Anda bisa berbicara bahasa manusia dengan jelas, Anda bisa memprogram AI.

Namun, tidak semua prompt diciptakan sama. Perbedaan antara prompt "tolong buatkan artikel" dengan prompt yang terstruktur bisa menghasilkan output yang sangat berbeda kualitasnya.

\section{Anatomi Prompt yang Baik}

Sebuah prompt yang efektif biasanya memiliki komponen-komponen berikut:

\begin{enumerate}
    \item \textbf{Role (Peran):} Siapa AI dalam konteks ini? (Misal: Guru, Editor, Ahli Gizi).
    \item \textbf{Context (Konteks):} Informasi latar belakang yang relevan.
    \item \textbf{Instruction (Instruksi):} Apa yang harus dilakukan AI secara spesifik.
    \item \textbf{Constraint (Batasan):} Apa yang tidak boleh dilakukan, atau format output yang diinginkan.
    \item \textbf{Few-Shot Examples (Contoh):} Contoh input dan output yang diharapkan.
\end{enumerate}

\begin{promptbox}[title=Contoh Prompt Terstruktur]
\textbf{Role:} Bertindaklah sebagai editor senior di koran nasional.

\textbf{Context:} Saya memiliki draf artikel tentang "Bahaya Plastik" yang ditulis dengan gaya bahasa terlalu santai.

\textbf{Instruction:} Tulis ulang artikel tersebut agar bernada formal, persuasif, dan mendesak.

\textbf{Constraint:} Jangan gunakan jargon teknis yang sulit dipahami orang awam. Maksimal 300 kata. Gunakan format bullet points untuk poin utama.

\textbf{Input:} [Artikel asli saya...]
\end{promptbox}

\section{Zero-Shot vs. Few-Shot Prompting}

\subsection{Zero-Shot}
Ini adalah saat Anda langsung meminta AI melakukan sesuatu tanpa memberikan contoh.
\begin{codeblock}
Prompt: Terjemahkan kalimat ini ke bahasa Jepang: "Saya suka makan nasi goreng."
\end{codeblock}

\subsection{Few-Shot (Pemberian Contoh)}
Teknik ini sangat ampuh untuk tugas yang kompleks atau membutuhkan format spesifik. Anda memberikan 1-3 contoh cara mengerjakannya.

\begin{promptbox}[title=Few-Shot Example]
Konversikan data mentah menjadi format JSON.

Input: Budi, 25 tahun, Jakarta
Output: {"nama": "Budi", "usia": 25, "kota": "Jakarta"}

Input: Siti, 30 tahun, Surabaya
Output: {"nama": "Siti", "usia": 30, "kota": "Surabaya"}

Input: Andi, 22 tahun, Bandung
Output:
\end{promptbox}

Dengan memberikan contoh, AI "mengerti" pola yang Anda inginkan tanpa perlu penjelasan panjang lebar.

\section{Chain of Thought (CoT)}

Chain of Thought adalah teknik meminta AI untuk "berpikir langkah demi langkah" sebelum memberikan jawaban akhir. Ini terbukti meningkatkan akurasi dalam soal matematika atau logika.

\begin{tipbox}
Tambahkan kalimat ajaib ini di akhir prompt Anda:
\textbf{"Let's think step by step."} (Mari berpikir langkah demi langkah).
\end{tipbox}

\begin{promptbox}[title=Perbandingan CoT]
\textbf{Tanpa CoT:}
Q: Jika saya punya 5 apel, makan 2, lalu beli 3 lagi, dan membagikan setengah dari totalnya ke teman, berapa sisa apel saya?
A: 3 apel. (Salah!)

\textbf{Dengan CoT:}
Q: Jika saya punya 5 apel, makan 2, lalu beli 3 lagi, dan membagikan setengah dari totalnya ke teman, berapa sisa apel saya? Mari berpikir langkah demi langkah.
A:
1. Awalnya ada 5 apel.
2. Dimakan 2, sisa 5 - 2 = 3.
3. Beli 3 lagi, jadi 3 + 3 = 6.
4. Membagikan setengahnya (6 / 2 = 3).
5. Sisa apel Anda adalah 3. (Benar!)
\end{promptbox}

\section{Teknik ReAct (Reason + Act)}

ReAct menggabungkan penalaran (Reasoning) dengan tindakan (Acting), biasanya digunakan ketika AI memiliki akses ke tools (seperti pencarian internet).

Pola: \textbf{Thought $\rightarrow$ Action $\rightarrow$ Observation}

\begin{itemize}
    \item \textbf{Thought:} Saya perlu mencari tahu siapa presiden Indonesia saat ini.
    \item \textbf{Action:} Search("Presiden Indonesia 2024").
    \item \textbf{Observation:} Hasil pencarian menunjukkan Joko Widodo (sampai Oktober) dan Prabowo Subianto (terpilih).
    \item \textbf{Thought:} Saya harus menjawab berdasarkan tanggal hari ini.
\end{itemize}

\section{Workshop: Memperbaiki Prompt}

\begin{exercisebox}
\textbf{Kasus:} Anda ingin membuat email penawaran jasa desain grafis.

\textbf{Prompt Buruk:}
"Buatkan email penawaran jasa desain."

\textbf{Tugas Anda:}
Ubah prompt di atas menjadi prompt yang menggunakan rumus \textbf{R-C-I-C} (Role, Context, Instruction, Constraint).

\vspace{2cm} % Space for writing
(Tulis jawaban Anda di sini...)
\end{exercisebox}

\section{Advanced Prompting Cheat Sheet}

\begin{center}
\begin{tabular}{|p{4cm}|p{8cm}|}
\hline
\textbf{Teknik} & \textbf{Pola Prompt} \\
\hline
\textbf{Self-Consistency} & "Generate 5 jawaban berbeda, lalu simpulkan jawaban yang paling sering muncul." \\
\hline
\textbf{Least-to-Most} & "Pecah masalah ini menjadi sub-masalah kecil. Selesaikan sub-masalah pertama, lalu gunakan hasilnya untuk yang kedua." \\
\hline
\textbf{Generated Knowledge} & "Tulis 5 fakta tentang topik ini, lalu gunakan fakta tersebut untuk menjawab pertanyaan saya." \\
\hline
\textbf{Meta-Prompting} & "Buatkan saya prompt terbaik untuk menyuruh AI melakukan X." \\
\hline
\end{tabular}
\end{center}
